\section{Why There Is Something Rather Than Nothing}
\label{sec:bootstrap}

The oldest question is why there is anything at all. The usual feeling is that nothing is the
simple, natural default, and that something is the strange thing that needs explaining. The
model takes the opposite view, and it is worth stating plainly, because it is what lets the
whole picture begin without a cause from outside. Nothing is the state that cannot hold.
Something is forced.

\paragraph{What nothing would have to mean.}
Be strict about the word. True nothing is not empty space, because empty space is already
something: it has size, it has directions, it has a clock. Strip all of that away. True nothing
would be the total absence of any distinction whatsoever: no space, no time, no things, no
states, no difference of any kind. It is not a dark room. It is no inside and no outside, with
no ``everything'' even to speak of.

\paragraph{Why nothing cannot hold.}
Three plain considerations point the same way. First, a state with no distinctions is a
balance point, not a resting place. A ball can in principle sit on the exact top of a
perfectly round hill, but the smallest nudge sends it rolling, and nothing pins it there. A
floor catches you. A hilltop does not. A perfectly uniform state is a hilltop. Second, ``no
difference'' is itself a difference. For there to be a pure uniform nothing, it would have to
be wholly the same all the way through, yet the moment it can be spoken of as one whole,
distinct from what it is not, a line has been drawn and a distinction made. The state meant to
contain no distinction cannot even be itself without making one. Third, a whole that is
everything has nothing else to relate to, so it can only relate to itself. And to exist, at
bottom, is to relate, since a thing that touches nothing and is touched by nothing is
indistinguishable from not being there. So the whole turns on itself, and the instant it does,
it splits into two roles, the side doing the relating and the side related to. That split is
the first difference. And there is no road back, because even erasing a distinction would be
an event, a change, a new difference between before and after, so the attempt to reach nothing
only makes more something. The door opens one way.

\paragraph{The bootstrap.}
Run it forward in the cleanest terms. Begin from undivided unity, the closest thing to nothing
that can even be spoken of, one whole with no parts. It cannot simply sit there, both because a
perfectly uniform state is a balance point that cannot hold, and because being one whole is
already a distinction. It can only relate to itself, there being nothing else. That
self-relation splits it in two, a noting end and a noted end, and this is the first
distinction. Everything basic is born in that one act. The two ends are the first two vibes,
the relation between them is the first note, the side each takes is the first tone, and the act
itself is the first beat, the first tick of time. Nothing was added from outside, because there
is no outside. The first state lifts itself by its own self-relation, since not doing so was
never a stable option. The beginning, on this reading, is not a lucky accident but a
necessity.

\begin{principle}[The forced beginning]
A state of no distinctions cannot hold. It is a balance point rather than a floor, it cannot be
named without making a distinction, and a whole with nothing else can only relate to itself,
which splits it. So a beginning is not something that needed an outside cause. It is the one
thing reality can do.
\end{principle}

\paragraph{The arrow, for free.}
Once the first distinction exists it is itself a new thing, and a new thing can be noted, and
noting it makes another distinction. Distinctions only accumulate, and what has happened cannot
be un-experienced. This one-way pile-up is the arrow of time. Time is not a track laid down in
advance for reality to move along. It is the plain fact that the heap of what has happened only
grows, the past large and fixed, the future open and being made. The direction comes for free
from the bootstrap, because distinctions can be added but not subtracted. This is the same
growth that, drawn out in space, became the ever-expanding mesh of the previous picture.

\paragraph{Where it heads, stated honestly.}
The reading the model commits to is eternal expansion: experience never ceases and never
exactly repeats, finite at every moment yet never exhausted, since there is always one more
distinction to make, with selves forming, coming into harmony, and dissolving up close while
the whole simply grows. A few alternatives are not yet ruled out, and each would change the
ending, so they are flagged as open. If genuine forgetting were possible, the story could
become an endless reshuffling of the same total, or even an exact return. If a changeless state
could somehow stay fully aware, a final still and conscious peace becomes thinkable, though that
reopens the original puzzle, since a perfectly still state is just the balance point that set
everything rolling. The necessity of a beginning does not depend on which fork is true. The
beginning is forced either way.
